
\subsubsection{7.1 Summary}

We designed a two-stage experimental framework to test the gradient competition hypothesis. First, we confirmed that spurious feature dominance exists: GradCAM attribution ratios exceed 1.0 from epoch 0, reaching 1.35 at initialization and peaking at 1.59 around epoch 13.

Second, we directly measured gradient norms for spurious-aligned versus minority samples. The results were unambiguous: the gradient ratio is 0.18, not the >1.5 predicted by gradient competition. Minority samples produce 5.$7\times$ higher gradient norms than spurious-aligned samples---a complete inversion of the hypothesis.

Our main findings are:

\begin{enumerate}
\item \textbf{Empirical falsification} of the gradient competition hypothesis with high confidence (approximately $8\times$ inversion, consistent across 3 seeds)
\end{enumerate}

\begin{enumerate}
\item \textbf{Confirmation and characterization} of simplicity bias dynamics, showing spurious dominance from initialization
\end{enumerate}

\begin{enumerate}
\item \textbf{Mechanistic reframing} proposing that spurious dominance arises from convergence speed---simpler patterns achieve low loss faster, not louder gradient signals
\end{enumerate}

\subsubsection{7.2 Future Directions}

\textbf{Testing the convergence speed hypothesis.} Direct measurement of per-group loss trajectories would confirm whether spurious-aligned samples reach low loss before minority samples.

\textbf{Connecting to loss landscape geometry.} Hessian eigenvalue analysis at peak spurious dominance versus convergence could reveal the geometric structure underlying our findings.

\textbf{Revisiting SAM through the convergence lens.} Our mechanistic reframing suggests SAM's effectiveness may stem from slowing convergence in simple basins.

\textbf{Multi-dataset replication.} CelebA, Colored MNIST, and CivilComments would test generalization of these findings.

\subsubsection{7.3 Closing Remarks}

The gradient competition hypothesis offered an intuitive explanation for spurious feature learning: simpler features receive stronger signals and thus dominate. Our measurements reveal this intuition is empirically incorrect. Spurious features dominate not because they produce louder gradients, but because they converge first.
